\documentclass[10pt,letterpaper]{article}
\usepackage[utf8]{inputenc}
\usepackage[intlimits]{amsmath}
\usepackage{amsfonts}
\usepackage{amssymb}
\usepackage{ragged2e}
\usepackage[letterpaper, margin=0.5in]{geometry}
\usepackage{graphicx}
\usepackage{cancel}
\usepackage{mathtools}
\usepackage{tabularx}
\usepackage{arydshln}
\usepackage{tensor}
\usepackage{array}
\usepackage{xcolor}
\usepackage[boxed]{algorithm}
\usepackage[noend]{algpseudocode}
\usepackage{listings}
\usepackage{textcomp}
% \usepackage[pdf,tmpdir,singlefile]{graphviz}
\usepackage{mathrsfs}
\usepackage{bbm}
\usepackage{tikz}
\usepackage{tikz-cd}
\usepackage{enumitem}
\usepackage{arydshln}
\usepackage{relsize}
\usepackage{multirow,multicol}
\usepackage{scalerel}
\usepackage{upgreek}
\usepackage{ifthen}

\usetikzlibrary{bayesnet}
\setlist{noitemsep}

%%%%%%%%%%%%%%%%%%%%%%%%%%%%%
% Formatting commands
%%%%%%%%%%%%%%%%%%%%%%%%%%%%%
\newcommand{\n}{\hfill\break}
\newcommand{\nn}{\vspace{0.5\baselineskip}\n}
\newcommand{\up}{\vspace{-\baselineskip}}
\newcommand{\hangblock}[2]{\par\noindent\settowidth{\hangindent}{\textbf{#1: }}\textbf{#1: }\nolinebreak#2}
\newcommand{\lemma}[1]{\hangblock{Lemma}{#1}}
\newcommand{\defn}[1]{\hangblock{Defn}{#1}}
\newcommand{\thm}[1]{\hangblock{Thm}{#1}}
\newcommand{\cor}[1]{\hangblock{Cor}{#1}}
\newcommand{\prop}[1]{\hangblock{Prop}{#1}}
\newcommand{\ex}[1]{\hangblock{Ex}{#1}}
\newcommand{\exer}[1]{\hangblock{Exer}{#1}}
\newcommand{\fact}[1]{\hangblock{Fact}{#1}}
\newcommand{\remark}[1]{\hangblock{Remark}{#1}}
\newcommand{\proven}{\;$\square$\n}
\newcommand{\problem}[1]{\par\noindent{\nolinebreak#1}\n}
\newcommand{\problempart}[2]{\par\noindent\indent{}\settowidth{\hangindent}{\textbf{(#1)} \indent{}}\textbf{(#1) }\nolinebreak#2\n}
\newcommand{\ptxt}[1]{\textrm{\textnormal{#1}}}
\newcommand{\inlineeq}[1]{\centerline{$\displaystyle #1$}}
\newcommand{\pageline}{\up\par\noindent\rule{\textwidth}{0.1pt}}

%%%%%%%%%%%%%%%%%%%%%%%%%%%%%
% Math commands
%%%%%%%%%%%%%%%%%%%%%%%%%%%%%
% Set Theory
\newcommand{\card}[1]{\left|#1\right|}
\newcommand{\set}[1]{\left\{#1\right\}}
\newcommand{\setmid}{\;\middle|\;}
\newcommand{\ps}[1]{\mathcal{P}\left(#1\right)}
\newcommand{\pfinite}[1]{\mathcal{P}^{\ptxt{finite}}\left(#1\right)}
\newcommand{\naturals}{\mathbb{N}}
\newcommand{\N}{\naturals}
\newcommand{\integers}{\mathbb{Z}}
\newcommand{\Z}{\integers}
\newcommand{\rationals}{\mathbb{Q}}
\newcommand{\Q}{\rationals}
\newcommand{\reals}{\mathbb{R}}
\newcommand{\R}{\reals}
\newcommand{\complex}{\mathbb{C}}
\newcommand{\C}{\complex}
\newcommand{\halfPlane}{\mathbb{H}}
\let\HSym\H
\let\H\relax
\newcommand{\H}{\halfPlane}
\newcommand{\comp}{^{\complement}}
\DeclareMathOperator{\Hom}{Hom}
\newcommand{\Ind}{\mathbbm{1}}
\newcommand{\cut}{\setminus}
\DeclareMathOperator{\elem}{elem}

% Differentiable Manifolds
\newcommand{\RP}{\mathbb{RP}}
\newcommand{\CP}{\mathbb{RP}}
\newcommand{\osubset}{\overset{\mathclap{\scalebox{0.5}{\ptxt{open}}}}{\subset}}
\newcommand{\osubseteq}{\overset{\mathclap{\scalebox{0.5}{\ptxt{open}}}}{\subseteq}}
\newcommand{\osupset}{\overset{\mathclap{\scalebox{0.5}{\ptxt{open}}}}{\supset}}
\newcommand{\osupseteq}{\overset{\mathclap{\scalebox{0.5}{\ptxt{open}}}}{\supseteq}}
\newcommand{\pdat}[3]{\left.\pd{#1}{#2}\right|_{#3}}
\DeclareMathOperator{\so}{so}
\DeclareMathOperator{\codim}{codim}
\DeclareMathOperator{\Diff}{Diff}
\let\dSym\d
\let\d\relax
\newcommand{\d}{\partial}
\DeclareMathOperator{\gl}{gl}
\DeclareMathOperator{\Ad}{Ad}
\newcommand{\comm}[1]{\left[#1\right]}
\DeclareMathOperator{\Fix}{Fix}

% Graph Theory
\let\deg\relax
\DeclareMathOperator{\deg}{deg}
\newcommand{\degp}{\ptxt{deg}^{+}}
\newcommand{\degn}{\ptxt{deg}^{-}}
\newcommand{\precdot}{\mathrel{\ooalign{$\prec$\cr\hidewidth\hbox{$\cdot\mkern0.5mu$}\cr}}}
\newcommand{\succdot}{\mathrel{\ooalign{$\cdot\mkern0.5mu$\cr\hidewidth\hbox{$\succ$}\cr\phantom{$\succ$}}}}
\DeclareMathOperator{\cl}{cl}
\DeclareMathOperator{\affdim}{affdim}

% Probability
\newcommand{\parSymbol}{\P}
\newcommand{\Prob}{\mathbb{P}}
\renewcommand{\P}{\Prob}
\newcommand{\Avg}{\mathbb{E}}
\newcommand{\E}{\Avg}
\DeclareMathOperator{\Var}{Var}
\DeclareMathOperator{\Cov}{Cov}
\DeclareMathOperator{\cov}{cov}
\DeclareMathOperator{\Unif}{Unif}
\DeclareMathOperator{\Binom}{Binom}
\newcommand{\CI}{\mathrel{\text{\scalebox{1.07}{$\perp\mkern-10mu\perp$}}}}
\DeclareMathOperator{\Ber}{Ber}
\DeclareMathOperator{\Bin}{Bin}
\DeclareMathOperator{\Geom}{Geom}
\DeclareMathOperator{\Poisson}{Poisson}
\DeclareMathOperator{\Exp}{Exp}
\DeclareMathOperator{\Beta}{Beta}
\DeclareMathOperator{\Normal}{N}
\DeclareMathOperator{\DistGamma}{Gamma}

% Standard Math
\newcommand{\inv}{^{-1}}
\newcommand{\abs}[1]{\left|#1\right|}
\newcommand{\ceil}[1]{\left\lceil{}#1\right\rceil{}}
\newcommand{\floor}[1]{\left\lfloor{}#1\right\rfloor{}}
\newcommand{\conj}[1]{\overline{#1}}
\newcommand{\of}{\circ}
\newcommand{\tri}{\triangle}
\newcommand{\inj}{\hookrightarrow}
\newcommand{\surj}{\twoheadrightarrow}
\newcommand{\ndiv}{\nmid}
\renewcommand{\epsilon}{\varepsilon}
\newcommand{\divides}{\mid}
\newcommand{\ndivides}{\nmid}
\DeclareMathOperator{\lcm}{lcm}
\DeclareMathOperator{\sgn}{sgn}
\newcommand{\map}[4]{\!\!\!\begin{array}[t]{rcl}#1 & \!\!\!\!\to & \!\!\!\!#2\\ {}#3 & \!\!\!\!\mapsto & \!\!\!\!#4\end{array}}
\newcommand{\bigsum}[2]{\smashoperator[lr]{\sum_{\scalebox{#1}{$#2$}}}}
\DeclareMathOperator{\gcf}{gcf}
\newcommand{\restr}[1]{\left.#1\right|}

% Linear Algebra
\newcommand{\Id}{\textrm{\textnormal{Id}}}
\newcommand{\im}{\textrm{\textnormal{im}}}
\newcommand{\norm}[1]{\abs{\abs{#1}}}
\newcommand{\tpose}{^{T}\!}
\newcommand{\iprod}[1]{\left<#1\right>}
\newcommand{\giprod}{\iprod{\;\,,\;}}
\DeclareMathOperator{\tr}{tr}
\DeclareMathOperator{\trace}{tr}
\newcommand{\chgBasMat}[3]{\!\!\tensor*[_{#1}]{\left[#2\right]}{_{#3}}}
\newcommand{\vecBas}[2]{\tensor*[]{\left[#1\right]}{_{#2}}}
\DeclareMathOperator{\GL}{GL}
\DeclareMathOperator{\Mat}{Mat}
\DeclareMathOperator{\vspan}{span}
\DeclareMathOperator{\rank}{rank}
\newcommand{\V}[1]{\vec{#1}}
\DeclareMathOperator{\proj}{proj}
\DeclareMathOperator{\compProj}{comp}
\DeclareMathOperator{\row}{row}
\newcommand{\smallPMatrix}[1]{\paren{\begin{smallmatrix}#1\end{smallmatrix}}}
\newcommand{\smallBMatrix}[1]{\brack{\begin{smallmatrix}#1\end{smallmatrix}}}
\newcommand{\pmat}[1]{\begin{pmatrix}#1\end{pmatrix}}
\newcommand{\bmat}[1]{\begin{bmatrix}#1\end{bmatrix}}
\newcommand{\dual}{^{*}}
\newcommand{\pinv}{^{\dagger}}
\newcommand{\horizontalMatrixLine}{\ptxt{\rotatebox[origin=c]{-90}{$|$}}}
\DeclareMathOperator{\range}{range}
\DeclareMathOperator{\Symm}{Symm}
\DeclareMathOperator{\SU}{SU}
\DeclareMathOperator{\U}{U}

% Multilinear Algebra
\let\Lsym\L
\let\L\relax
\DeclareMathOperator{\L}{\mathscr{L}}
\DeclareMathOperator{\A}{\mathcal{A}}
\DeclareMathOperator{\Alt}{Alt}
\DeclareMathOperator{\Sym}{Sym}
\newcommand{\ot}{\otimes}
\newcommand{\ox}{\otimes}
\DeclareMathOperator{\asc}{asc}
\DeclareMathOperator{\asSet}{set}
\DeclareMathOperator{\sort}{sort}
\DeclareMathOperator{\ringA}{\mathring{A}}
\DeclareMathOperator{\Sh}{Sh}
\DeclareMathOperator{\Bil}{Bil}

% Topology
\newcommand{\closure}[1]{\overline{#1}}
\newcommand{\uball}{\mathcal{U}}
\DeclareMathOperator{\Int}{Int}
\DeclareMathOperator{\Ext}{Ext}
\DeclareMathOperator{\Bd}{Bd}
\DeclareMathOperator{\rInt}{rInt}
\DeclareMathOperator{\ch}{ch}
\DeclareMathOperator{\ah}{ah}
\newcommand{\LargerTau}{\mathlarger{\mathlarger{\mathlarger{\mathlarger{\tau}}}}}
\newcommand{\Tau}{\mathcal{T}}

% Analysis
\DeclareMathOperator{\Graph}{Graph}
\DeclareMathOperator{\epi}{epi}
\DeclareMathOperator{\hypo}{hypo}
\DeclareMathOperator{\supp}{supp}
\newcommand{\lint}[2]{\underset{#1}{\overset{#2}{{\color{black}\underline{{\color{white}\overline{{\color{black}\int}}\color{black}}}}}}}
\newcommand{\uint}[2]{\underset{#1}{\overset{#2}{{\color{white}\underline{{\color{black}\overline{{\color{black}\int}}\color{black}}}}}}}
\newcommand{\alignint}[2]{\underset{#1}{\overset{#2}{{\color{white}\underline{{\color{white}\overline{{\color{black}\int}}\color{black}}}}}}}
\newcommand{\extint}{\ptxt{ext}\int}
\newcommand{\extalignint}[2]{\ptxt{ext}\alignint{#1}{#2}}
\newcommand{\conv}{\ast}
\newcommand{\pd}[2]{\frac{\partial{}#1}{\partial{}#2}}
\newcommand{\del}{\nabla}
\DeclareMathOperator{\grad}{grad}
\DeclareMathOperator{\curl}{curl}
\let\div\relax
\DeclareMathOperator{\div}{div}
\DeclareMathOperator{\vol}{vol}

% Complex Analysis
\let\Re\relax
\DeclareMathOperator{\Re}{Re}
\let\Im\relax
\DeclareMathOperator{\Im}{Im}
\DeclareMathOperator{\Res}{Res}

% Abstract Algebra
\DeclareMathOperator{\ord}{ord}
\newcommand{\generated}[1]{\left<#1\right>}
\newcommand{\cycle}[1]{\smallPMatrix{#1}}
\newcommand{\id}{\ptxt{id}}
\newcommand{\iso}{\cong}
\DeclareMathOperator{\Aut}{Aut}
\DeclareMathOperator{\SL}{SL}
\DeclareMathOperator{\op}{op}
\newcommand{\isom}[4]{\!\!\!\begin{array}[t]{rcl}#1 & \!\!\!\!\overset{\sim}{\to} & \!\!\!\!#2\\ #3 & \!\!\!\!\mapsto & \!\!\!\!#4\end{array}}
\newcommand{\F}{\mathbb{F}}
\newcommand{\acts}{\;\reflectbox{\rotatebox[origin=c]{-90}{$\circlearrowright$}}\;}
\newcommand{\disjunion}{\mathrel{\text{\scalebox{1.07}{$\perp\mkern-10mu\perp$}}}}
\DeclareMathOperator{\SO}{SO}
\DeclareMathOperator{\stab}{stab}
\DeclareMathOperator{\nullity}{nullity}
\DeclareMathOperator{\Perm}{Perm}
\DeclareMathOperator{\nsubgp}{\vartriangleleft}
\DeclareMathOperator{\notnsubgp}{\ntriangleleft}
\newcommand{\presentation}[2]{\left<#1\;\middle|\;#2\right>}
\DeclareMathOperator{\Char}{char}
\DeclareMathOperator{\fchar}{char}
\DeclareMathOperator{\triv}{triv}
\DeclareMathOperator{\reg}{reg}
\DeclareMathOperator{\std}{std}
\DeclareMathOperator{\Func}{Func}
\DeclareMathOperator{\End}{End}

% Convex Optimization
\let\sectionSymbol\S
\let\S\relax
\newcommand{\S}{\mathbb{S}}
\DeclareMathOperator{\dist}{dist}
\DeclareMathOperator{\dom}{dom}
\DeclareMathOperator{\diag}{diag}
\DeclareMathOperator{\ones}{\mathbbm{1}}
\newcommand{\minimizeOver}[3]{\begin{array}{rl}\underset{#1}{\ptxt{minimize}} & #2\\ \ptxt{subject to} & #3\end{array}}
\newcommand{\maximizeOver}[3]{\begin{array}{rl}\underset{#1}{\ptxt{maximize}} & #2\\ \ptxt{subject to} & #3\end{array}}
\newcommand{\minimizationProblem}[2]{\minimizeOver{}{#1}{#2}}
\newcommand{\maximizationProblem}[2]{\maximizeOver{}{#1}{#2}}
\newcommand{\minimizeOverUnconstrained}[2]{\begin{array}{rl}\underset{#1}{\ptxt{minimize}} & #2\end{array}}
\newcommand{\maximizeOverUnconstrained}[2]{\begin{array}{rl}\underset{#1}{\ptxt{maximize}} & #2\end{array}}
\newcommand{\minimizationUnconstrained}[1]{\minimizeOverUnconstrained{}{#1}}
\newcommand{\maximizationUnconstrained}[1]{\maximizeOverUnconstrained{}{#1}}
\DeclareMathOperator{\argmin}{argmin}
\DeclareMathOperator{\argmax}{argmax}

% Proofs
\newcommand{\st}{s.t.}
\newcommand{\unique}{!}
\newcommand{\iffdef}{\overset{\ptxt{def}}{\Leftrightarrow}}
\newcommand{\eqVertical}{\rotatebox[origin=c]{90}{=}}
\newcommand{\mapsfrom}{\mathrel{\reflectbox{\ensuremath{\mapsto}}}}
\newcommand{\mapsdown}{\rotatebox[origin=c]{-90}{$\mapsto$}\mkern2mu}
\newcommand{\mapsup}{\rotatebox[origin=c]{90}{$\mapsto$}\mkern2mu}
\newcommand{\from}{\!\mathrel{\reflectbox{\ensuremath{\to}}}}
\newcommand{\labeledeq}[1]{\overset{\mathclap{\ptxt{#1}}}{=}}
\newcommand{\eqdef}{\labeledeq{def}}

% Brackets
\newcommand{\paren}[1]{\left(#1\right)}
\renewcommand{\brack}[1]{\left[#1\right]}
\renewcommand{\brace}[1]{\left\{#1\right\}}
\newcommand{\ang}[1]{\left<#1\right>}

% Algorithms
\algrenewcommand{\algorithmiccomment}[1]{\hskip 1em \texttt{// #1}}
\algrenewcommand\algorithmicrequire{\textbf{Input:}}
\algrenewcommand\algorithmicensure{\textbf{Output:}}
\newcommand{\algP}{\ptxt{\textbf{P}}}
\newcommand{\algNP}{\ptxt{\textbf{NP}}}
\newcommand{\algNPC}{\ptxt{\textbf{NP-Complete}}}
\newcommand{\algNPH}{\ptxt{\textbf{NP-Hard}}}
\newcommand{\algEXP}{\ptxt{\textbf{EXP}}}
\DeclareMathOperator{\fl}{fl}

%%%%%%%%%%%%%%%%%%%%%%%%%%%%%
% Other commands and shorthand
%%%%%%%%%%%%%%%%%%%%%%%%%%%%%
\newcommand{\flag}[1]{\textbf{\textcolor{red}{#1}}}
\let\uSym\u
\let\u\relax
\newcommand{\u}[1]{\underline{#1}}
\let\bSym\b
\let\b\relax
\newcommand{\b}[1]{\textbf{#1}}
\let\iSym\i
\let\i\relax
\newcommand{\i}[1]{\textit{#1}}
\let\scSym\sc
\let\sc\relax
\newcommand{\sc}[1]{\textsc{#1}}
\newcommand{\mf}[1]{\mathfrak{#1}}
\newcommand{\mc}[1]{\mathcal{#1}}

%%%%%%%%%%%%%%%%%%%%%%%%%%%%%%%%%%%%%%%
% Make l's curvy in math environments %
%%%%%%%%%%%%%%%%%%%%%%%%%%%%%%%%%%%%%%%
\mathcode`l="8000
\begingroup
\makeatletter
\lccode`\~=`\l
\DeclareMathSymbol{\lsb@l}{\mathalpha}{letters}{`l}
\lowercase{\gdef~{\ifnum\the\mathgroup=\m@ne \ell \else \lsb@l \fi}}%
\endgroup

%%%%%%%%%%%%%%%%%%%%%%%%%
% Fix \vdots and \ddots %
%%%%%%%%%%%%%%%%%%%%%%%%%
\usepackage{letltxmacro}
\LetLtxMacro\orgvdots\vdots
\LetLtxMacro\orgddots\ddots

\makeatletter
\DeclareRobustCommand\vdots{%
	\mathpalette\@vdots{}%
}
\newcommand*{\@vdots}[2]{%
	% #1: math style
	% #2: unused
	\sbox0{$#1\cdotp\cdotp\cdotp\m@th$}%
	\sbox2{$#1.\m@th$}%
	\vbox{%
		\dimen@=\wd0 %
		\advance\dimen@ -3\ht2 %
		\kern.5\dimen@
		% remove side bearings
		\dimen@=\wd2 %
		\advance\dimen@ -\ht2 %
		\dimen2=\wd0 %
		\advance\dimen2 -\dimen@
		\vbox to \dimen2{%
			\offinterlineskip
			\copy2 \vfill\copy2 \vfill\copy2 %
		}%
	}%
}
\DeclareRobustCommand\ddots{%
	\mathinner{%
		\mathpalette\@ddots{}%
		\mkern\thinmuskip
	}%
}
\newcommand*{\@ddots}[2]{%
	% #1: math style
	% #2: unused
	\sbox0{$#1\cdotp\cdotp\cdotp\m@th$}%
	\sbox2{$#1.\m@th$}%
	\vbox{%
		\dimen@=\wd0 %
		\advance\dimen@ -3\ht2 %
		\kern.5\dimen@
		% remove side bearings
		\dimen@=\wd2 %
		\advance\dimen@ -\ht2 %
		\dimen2=\wd0 %
		\advance\dimen2 -\dimen@
		\vbox to \dimen2{%
			\offinterlineskip
			\hbox{$#1\mathpunct{.}\m@th$}%
			\vfill
			\hbox{$#1\mathpunct{\kern\wd2}\mathpunct{.}\m@th$}%
			\vfill
			\hbox{$#1\mathpunct{\kern\wd2}\mathpunct{\kern\wd2}\mathpunct{.}\m@th$}%
		}%
	}%
}
\makeatother

\newcommand{\B}{
	\begin{tikzpicture}
	\filldraw [fill=red, draw=black] (0, 0) rectangle (0.37, 0.45);
	\draw [line width=0.5mm, white ] (0.1,0.08) -- (0.1,0.38);
	\draw[line width=0.5mm, white ] (0.1, 0.35) .. controls (0.2, 0.35) and (0.4, 0.2625) .. (0.1, 0.225);
	\draw[line width=0.5mm, white ] (0.1, 0.225) .. controls (0.2, 0.225) and (0.4, 0.1625) .. (0.1, 0.1);
	\end{tikzpicture}
}

% Allow custom symbols as arrows in tikz-cd
\tikzset{
  symbol/.style={
    draw=none,
    every to/.append style={
      edge node={node [sloped, allow upside down, auto=false]{$#1$}}}
  }
}
% Example: \arrow[r,symbol=\cong]
% https://tex.stackexchange.com/questions/394154/how-to-include-inclusion-subgroup-relationship-in-tikz-cd-diagram

\author{Thomas Cohn}
\title{Stats 426 Lecture 6}
\date{2/8/21} % Can also use \today

\begin{document}
\maketitle
\setlength\RaggedRightParindent{\parindent}
\RaggedRight

\section*{Functions of Random Variables}

\par\noindent
Suppose $(Y_{1},\ldots,Y_{m})=g(X_{1},\ldots,X_{n})$. We want to find the joint pdf of $(Y_{1},\ldots,Y_{m})$, using the pdf of $(X_{1},\ldots,X_{n})$, denoted $f_{X_{1},\ldots,X_{n}}(x_{1},\ldots,x_{n})$. In the discrete case, we have
\[
	f_{Y_{1},\ldots,Y_{m}}(y_{1},\ldots,y_{m})=\sum_{\mathclap{\substack{(x_{1},\ldots,x_{n}):\\ g(x_{1},\ldots,x_{n})=(y_{1},\ldots,y_{m})}}}f_{X_{1},\ldots,X_{n}}(x_{1},\ldots,x_{n})
\]
In the continuous case, we can find the cdf, and then differentiate to get the pdf. Alternatively, if $m=n$, we can use the change of variables rule:
\[
	dy_{1}\cdots{}dy_{n}=\abs{\det{}J_{g}(x_{1},\ldots,x_{n})}\,dx_{1}\cdots{}dx_{n}
\]
where $J_{g}(x_{1},\ldots,x_{n})$ is the Jacobian, the matrix of partial derivatives
\[
	J_{g}(x_{1},\ldots,x_{n})=\begin{bmatrix}
		\pd{g_{1}}{x_{1}}(x_{1},\ldots,x_{n}) & \cdots & \pd{g_{1}}{x_{n}}(x_{1},\ldots,x_{n})\\
		\vdots & & \vdots\\
		\pd{g_{n}}{x_{1}}(x_{1},\ldots,x_{n}) & \cdots & \pd{g_{n}}{x_{n}}(x_{1},\ldots,x_{n})
	\end{bmatrix}
\]

\par\noindent
Suppose $X$ and $Y$ are two i.i.d. random variables, and $Z=X+Y$. In the discrete case, with probability mass functions $p,p_{X},p_{Y}$, then
\[
	p_{Z}(z)=P(X+Y=z)=\sum_{\mathclap{(x,y):x+y=z}}p(x,y)=\sum_{x\in{}X}p(x,z-x)\labeledeq{*}\sum_{x\in{}X}p_{X}(x)p_{Y}(z-x)
\]
*only if $X$ and $Y$ are independent. In the continuous case, with probability mass functions $f,f_{X},f_{Y}$, we have
\[
	f_{Z}(z)=\int_{X}f(x,z-x)\,dx\labeledeq{*}\int_{X}f_{X}(x)f_{Y}(z-x)\,dx
\]
*only if $X$ and $Y$ are independent.\n

\ex{
	Let $T_{1},T_{2}\sim\Exp(\lambda)$ be independent, and $S=T_{1}+T_{2}$. Then
	\[
		f_{S}(s)=\int_{\R}f_{T_{1}}(x)f_{T_{2}}(s-x)\,dx=\int_{0}^{s}\lambda{}e^{-\lambda{}x}\lambda{}e^{-\lambda(s-x)}\,dx=\int_{0}^{s}\lambda^{2}e^{-\lambda{}s}\,dx=s\lambda^{2}e^{-\lambda{}s}
	\]
	Thus, $S\sim\DistGamma(\lambda,z)$ for some constant $z$, that depends on $\lambda$.\n
}

\ex{
	Let $S^{[n]}=\sum_{i=1}^{n}T_{i}$, where the $T_{i}\sim\Exp(\lambda)$ are i.i.d. We just showed that $S^{[2]}=S^{[1]}+T$, where $T\sim\DistGamma(\lambda,z)$, and $T\perp{}S^{[1]}$. Suppose for some $k$, $S^{[k]}=S^{[k-1]}+T\sim\DistGamma(\lambda,k)$. Then
	\[
		\int_{0}^{s}\frac{\lambda^{k}}{\Gamma(k)}x^{k-1}e^{-\lambda{}k}\cdot\lambda{}e^{-\lambda(s-x)}\,dx=\int_{0}^{s}\frac{\lambda^{k+1}}{\Gamma(k)}x^{k-1}e^{-\lambda{}s}\,dx=\brack{\frac{\lambda^{k+1}}{k\Gamma(k)}e^{-\lambda{}s}x^{k}}_{x=0}^{x=s}=\frac{\lambda^{k+1}}{\Gamma(k+1)}s^{k}e^{-\lambda{}s}
	\]
	Thus, $S^{[k+1]}=S^{[k]}+T\sim\DistGamma(\lambda,k+1)$.\n
}

\newpage
\subsection*{Transformation}

\par\noindent
Consider $(Y_{1},Y_{2})=g(X_{1},X_{2})=(g_{1}(X_{1},X_{2}),g_{2}(X_{1},X_{2}))$, where $g$ is differentiable and invertible. Also, suppose the Jacobian is nonzero, $\forall(x_{1},x_{2})\in\supp{}f_{X_{1},X_{2}}$, and that $g$ maps $\supp{}f_{X_{1},X_{2}}$ to $\supp{}f_{Y_{1},Y_{2}}$. Then the joint probability density function of $(Y_{1},Y_{2})$ is
\[
	f_{Y_{1},Y_{2}}(y_{1},y_{2})=f_{X_{1},X_{2}}(h_{1}(y_{1},y_{2}),h_{2}(y_{1},y_{2}))J\inv(h_{1}(y_{1},y_{2}),h_{2}(y_{1},y_{2}))
\]
for $(y_{1},y_{2})\in\supp{}f_{Y_{1},Y_{2}}$ and $0$ otherwise, where $x_{1}=h_{1}(y_{1},y_{2})$, $x_{2}=h_{2}(y_{1},y_{2})$, and
\[
	J(x_{1},x_{2})=\det \begin{bmatrix}
		\pd{g_{1}}{x_{1}} & \pd{g_{1}}{x_{2}}\\
		\pd{g_{2}}{x_{1}} & \pd{g_{2}}{x_{2}}
	\end{bmatrix}
\]

\ex{
	The pdf of a ratio. Suppose $X$ and $Y$ are continuous random variables with joint pdf $f$, and maginal pdfs $f_{X}$ and $f_{Y}$. Let $Z=Y/X$. We can find the marginal pdf of $Z$ in two ways.\nn
	1: Change $(X,Y)$ to $(X,XZ)$. Then the Jacobian is
	\[
		\begin{bmatrix}
			\pd{X}{X} & \pd{X}{Z}\\
			\pd{Y}{X} & \pd{Y}{Z}
		\end{bmatrix}=\begin{bmatrix}
			1 & 0\\
			Z & X
		\end{bmatrix}
	\]
	So
	\[
		f_{X,Z}(x,z)=f_{X,Y}(x,xz)\abs{x}\qquad{}f_{Z}(z)=\int_{-\infty}^{\infty}f_{X,Y}(x,xz)\abs{x}\,dx
	\]
	2: Change $(X,Y/X)$ to $(X,Z)$. Then the Jacobian is
	\[
		\begin{bmatrix}
			\pd{X}{X} & \pd{X}{Y}\\
			\pd{Z}{X} & \pd{Z}{Y}
		\end{bmatrix}=\begin{bmatrix}
			1 & 0\\
			-YX^{2} & 1/X
		\end{bmatrix}=\begin{bmatrix}
			1 & 0\\
			-Z/X & 1/X
		\end{bmatrix}
	\]
	So
	\[
		f_{X,Z}(x,z)=f_{X,Y}(x,xz)\abs{\frac{1}{x}}\inv=f_{X,Y}(x,xz)\abs{x}
	\]
	We get the same thing -- both methods work!\n
}

\par\noindent
In general, if $X$ and $Y$ are independent, then
\[
	f_{Z}(z)=\int_{-\infty}^{\infty}\abs{x}f(x,xz)\,dx=\int_{-\infty}^{\infty}\abs{x}f_{X}(x)f_{Y}(xz)\,dx
\]

\ex{
	If $X,Y\sim\Normal(0,1)$ are i.i.d., show that the pdf of $Z=Y/X$ is $\frac{1}{\pi(1+z^{2})}$. Well,
	\begin{align*}
		f_{Z}(z) & =\int_{-\infty}^{\infty}\frac{\abs{x}}{\sqrt{2\pi}}e^{-x^{2}/2}\cdot\frac{1}{\sqrt{2\pi}}e^{-x^{2}z^{2}/2}\,dx\\
		& =\int_{-\infty}^{0}\frac{-x}{2\pi}e^{-(1+z^{2})x^{2}/2}\,dx+\int_{0}^{\infty}\frac{x}{2\pi}e^{-(1+z^{2})x^{2}/2}\,dx\qquad\begin{smallmatrix}s=x^{2}\\ ds=2x\,dx\end{smallmatrix}\\
		& =\int_{0}^{\infty}\frac{1}{4\pi}e^{-(1+z^{2})z/2}\,dx+\int_{0}^{\infty}\frac{1}{4\pi}e^{-(1+z^{2})s/2}\,ds\\
		& =\frac{1}{2\pi}\int_{0}^{\infty}e^{-(1+z^{2})s/2}\,ds\\
		& =\frac{1}{2\pi}\brack{\frac{-2}{1+z^{2}}e^{-(1+z^{2})s/2}}_{s=0}^{s=\infty}\\
		& =\frac{1}{\pi(1+z^{2})}
	\end{align*}
	\proven
}

\newpage
\ex{
	Suppose $X$ and $Y$ are i.i.d. standard normal random variables. Consider the polar coordinate system $X=R\cos\Theta$, $Y=R\sin\Theta$, for $R\ge{}0$ and $-\pi<\Theta<\pi$. Find the joint probability density function of $(R,\Theta)$. Well, we change $(X,Y)$ to $(R\cos\Theta,R\sin\Theta)$. So the Jacobian is
	\[
		\begin{bmatrix}
			\pd{X}{R} & \pd{X}{\Theta}\\
			\pd{Y}{R} & \pd{Y}{\Theta}
		\end{bmatrix}=\begin{bmatrix}
			\cos\Theta & -R\sin\Theta\\
			\sin\Theta & R\cos\Theta
		\end{bmatrix}
	\]
	So $\det{}J=R\cos^{2}\Theta-(-R\sin^{2}\Theta)=R\cos^{2}\Theta+R\sin^{2}\Theta=R$. Therefore,
	\[
		f_{R,\Theta}(r,\theta)=\frac{1}{2\pi}\abs{r}e^{-r^{2}/2}=\frac{r}{2\pi}e^{-r^{2}/2}\qquad\ptxt{for}\;r\ge{}0,-\pi<\theta<\pi
	\]
}

\exer{
	Use the other method -- change $(\sqrt{X^{2}+Y^{2}},\arctan(Y/X))$ to $(R,\Theta)$. Check that you get the same result.\n
}

\ex{
	Show that $R,\Theta$ are independent. Well,
	\begin{align*}
		& f_{R}(r)=\int_{-\pi}^{\pi}\frac{r}{2\pi}e^{-r^{2}/2}\,d\theta=re^{-r^{2}/2}\qquad\ptxt{for}\;r\ge{}0\\
		& f_{\Theta}{\theta}=\int_{0}^{\infty}\frac{r}{2\pi}e^{-r^{2}/2}\,dr=\cdots=\brack{\frac{-1}{2\pi}e^{-r^{2}/2}}_{r=0}^{r=\infty}=\frac{1}{2\pi}
	\end{align*}
	So
	\[
		f_{R}(r)f_{\Theta}(\theta)=\frac{r}{2\pi}e^{-r^{2}/2}=f_{R,\Theta}(r,\theta)\qquad\Rightarrow\qquad{}R\perp\Theta
	\]
	\proven
}

\subsection*{Expected Value}

\defn{
	Suppose $X$ is a discrete random variable, with pmf $p$. Then the \u{expected value} of $X$ is given by $\displaystyle{}E(X)=\sum_{x\in{}X}xp(x)$ if $\displaystyle\sum_{x\in{}X}\abs{x}p(x)<\infty$. If this sum is infinite, then the expected value is undefined.\n
}

\defn{
	Suppose $X$ is a continuous random variable, with pdf $f$. Then the \u{expected value} of $X$ is given by $\displaystyle{}E(X)=\int_{X}xp(x)\,dx$ if $\displaystyle\int_{X}\abs{x}f(x)\,dx<\infty$. If this integral is infinite, then the expected value is undefined.\n
}

\par\noindent
For constants $a,b$, $E(aX+b)=aE(X)+b$.\n

\exer{
	Derive the following:
	\begin{itemize}
		\item $X\sim\Ber(p)\rightsquigarrow{}E(X)=p$
		\item $X\sim\Bin(n,p)\rightsquigarrow{}E(X)=np$
		\item $X\sim\Geom(p)\rightsquigarrow{}E(X)=1/p$
		\item $X\sim\Poisson(\lambda)\rightsquigarrow{}E(X)=\lambda$
		\item $X\sim\Unif(a,b)\rightsquigarrow{}E(X)=(a+b)/2$
		\item $X\sim\Exp(\lambda)\rightsquigarrow{}E(X)=1/\lambda$
		\item $X\sim\DistGamma(\lambda,\alpha)\rightsquigarrow{}E(X)=\alpha/\lambda$
		\item $X\sim\Beta(a,b)\rightsquigarrow{}E(X)=a/(a+b)$
		\item $X\sim\Normal(\mu,\sigma^{2})\rightsquigarrow{}E(X)=\mu$
	\end{itemize}
}

\end{document}